\section{Test-Set Results and Error Analysis}\label{sec:test}

\subsection{Protocol}\label{sec:test-protocol}
The test split (1,892 images: 1,123 authentic, 769 tampered; 1,092 split groups, none shared with the
development set) was frozen when the splits were made (Section~\ref{sec:leakage}). It was used \emph{once} for the
headline evaluation (run~\#1) and afterwards only for post-hoc evaluations with the same frozen models (the
robustness and synthetic-forgery experiments of Section~\ref{sec:robustness}) and for the pre-registered, one-time
comparison with an ELA-CNN baseline (Section~\ref{sec:cnn}; new models, evaluated once). Nothing from these
experiments was fed back into any model or setting, and each use is recorded in the same run log.

All classifiers, calibrated models and localisers were frozen during development, and the run was guarded by dry
runs, an independent review, a preflight check and a fingerprinted audit record (Appendix~\ref{app:test}).
Confidence intervals are 95\,\% group-bootstrap intervals (1,000 resamples of test split groups).

\subsection{Classification}\label{sec:test-classification}
Table~\ref{tab:test-auc} and Fig.~\ref{fig:roc-test} show that the development estimates hold on unseen data.
Every test AUC is within 0.02 of its cross-validated development estimate (within 0.015 for all random-forest
rows), and every dev estimate falls inside the corresponding test CI (R forensic: dev 0.795, test
0.805\ci{0.779}{0.831}). There is no sign of overfitting to the development set, despite the feature tuning and
model selection done on it.

The central claim is confirmed on test (Table~\ref{tab:test-added}). Under the strict protocol R, the forensic
features beat all global shortcuts combined (0.805 vs 0.712). Added on top of them, they raise AUC by
+0.103\ci{0.076}{0.129}; the purely within-image features alone add +0.083\ci{0.048}{0.116}. Under B the gains
are smaller but still significant (+0.050; local only +0.022\ci{0.003}{0.040}). Protocol A again shows the leak:
on the released files, shortcuts alone reach test AUC 1.000 and forensic features add nothing. RF and SVM cannot
be told apart on the forensic features (McNemar, R: 119 vs 115 discordant, $p=0.84$; B: $p=0.88$).
AUC is the primary metric because the classes are imbalanced (59/41 authentic/tampered) and the operating point
of the deployed detector is set by calibration and abstention (Section~\ref{sec:test-deployed}); for completeness,
Table~\ref{tab:test-threshold} gives the threshold metrics of the forensic random forest at probability 0.5.



\begin{table*}[t]
\centering
\caption{Test ROC-AUC [95\,\% CI], with the development estimate (Section~\ref{sec:classification}) for
comparison.}\label{tab:test-auc}
\footnotesize
\begin{tabular}{lcccccc}
\toprule
Features & R: RF, test & R: RF, dev & R: SVM, test & B: RF, test & B: RF, dev & A: RF, test \\
\midrule
Shortcuts                  & 0.712\ci{0.682}{0.746} & 0.721 & 0.753 & 0.886 & 0.876 & \textbf{1.000} \\
Forensic, local only       & 0.796\ci{0.769}{0.821} & 0.781 & --    & 0.909 & 0.902 & 0.963 \\
\textbf{Forensic, local + global} & \textbf{0.805\ci{0.779}{0.831}} & 0.795 & 0.809 & \textbf{0.920} & 0.910 & 0.981 \\
Forensic + shortcuts       & 0.815\ci{0.789}{0.840} & 0.810 & 0.829 & 0.936 & 0.926 & 1.000 \\
\bottomrule
\end{tabular}

\vspace{3mm}
\caption{Test added value over the shortcuts (paired group bootstrap, $\Delta$AUC [95\,\% CI]).}\label{tab:test-added}
\begin{tabular}{lcccc}
\toprule
Contrast & R: RF & R: SVM & B: RF & A: RF \\
\midrule
(Forensic + shortcuts) $-$ shortcuts        & \textbf{+0.103\ci{0.076}{0.129}} & +0.077\ci{0.054}{0.099} & +0.050\ci{0.036}{0.063} & +0.000 \\
Forensic (local + global) $-$ shortcuts     & +0.093\ci{0.059}{0.125} & +0.056\ci{0.026}{0.085} & +0.033\ci{0.017}{0.049} & $-$0.019 \\
Forensic (local only) $-$ shortcuts         & \textbf{+0.083\ci{0.048}{0.116}} & -- & +0.022\ci{0.003}{0.040} & $-$0.037 \\
\bottomrule
\end{tabular}
\end{table*}

\begin{table}[t]
\centering
\caption{Test metrics at threshold 0.5 (forensic features, local + global, random forest; tampered is the positive
class).}\label{tab:test-threshold}
\footnotesize
\setlength{\tabcolsep}{4pt}
\begin{tabular}{lccccc}
\toprule
 & Accuracy & Precision & Recall & F1 & Balanced acc. \\
\midrule
R & 0.749 & 0.756 & 0.567 & 0.648 & 0.721\ci{0.695}{0.748} \\
B & 0.850 & 0.813 & 0.819 & 0.816 & 0.845\ci{0.826}{0.863} \\
\bottomrule
\end{tabular}
\end{table}

\begin{figure}[t]
\centering
\includegraphics[width=\columnwidth]{fig_roc_test_R.png}
\caption{ROC curves of the forensic random forest under protocol R: development out-of-fold (AUC 0.795) and
the one-time test evaluation (AUC 0.805). The two curves nearly coincide.}
\label{fig:roc-test}
\end{figure}

\subsection{The Deployed Detector}\label{sec:test-deployed}
The deployed detector (protocol R, calibrated, uncertain band $0.5\pm0.24$) behaves on test as it did in
development: it judges 51.8\,\% of the images (dev out-of-fold 51\,\%), with accuracy 0.858 (0.855) and balanced
accuracy 0.809 (0.810) on the judged images, and an expected calibration error of 0.019 (0.009; binned ECE rises
on smaller samples: test 1,892, dev about 10.7k). Table~\ref{tab:test-deployed} gives its verdicts. A ``tampered''
verdict is right 94\,\% of the time (222/235); an ``authentic'' verdict is right 83\,\% of the time (619/745),
because a missed forgery (126) is more common than a false alarm (13). Under B (band $\pm0.04$), 95.7\,\% of
images are judged, with accuracy 0.869 and balanced accuracy 0.853.

\begin{table}[t]
\centering
\caption{Verdicts of the deployed detector on the 1,892 test images (R, band $0.5\pm0.24$).}\label{tab:test-deployed}
\footnotesize
\begin{tabular}{lccc}
\toprule
True class & Authentic & Tampered & Uncertain \\
\midrule
Authentic (1,123) & 619 & 13  & 491 \\
Tampered (769)    & 126 & 222 & 421 \\
\bottomrule
\end{tabular}
\end{table}

\subsection{Localisation}\label{sec:test-localisation}
Table~\ref{tab:test-loc} reports the deployed localiser with frozen post-processing on the 766 test tampered
images with a valid mask. Under R, test localisation matches validation (F1 0.198 vs 0.185); under B it is
\emph{higher} than validation (0.308 vs 0.260, above the test CI); the cause was not investigated. The learned
fusion beats both the best single cue and the plain average; paired test differences in F1 are, under R,
+0.023\ci{0.009}{0.037} over keypoints and +0.029\ci{0.013}{0.045} over the average, and under B
+0.096\ci{0.076}{0.114} and +0.101\ci{0.080}{0.120}. Gating the map by the image verdict reduces the share of
authentic images showing a false region from 89\,\% to 43\,\% (R) and from 73\,\% to 9\,\% (B), at the cost of
tampered F1 falling from 0.198 to 0.189 (R) and from 0.308 to 0.267 (B).

\begin{table}[t]
\centering
\caption{Test localisation (766 tampered images with a valid mask; pixel F1 with 95\,\% CI). Baselines use
post-processing tuned on validation.}\label{tab:test-loc}
\footnotesize
\setlength{\tabcolsep}{4pt}
\begin{tabular}{lcc}
\toprule
 & R & B \\
\midrule
\textbf{Pixel F1} & \textbf{0.198\ci{0.179}{0.216}} & \textbf{0.308\ci{0.284}{0.332}} \\
IoU        & 0.128 & 0.210 \\
MCC        & 0.163 & 0.296 \\
Pixel AUC  & 0.714 & 0.829 \\
F1, mean of maps                  & 0.169 & 0.208 \\
F1, copy-move keypoints only      & 0.174 & 0.212 \\
F1, whole image                   & 0.135 & 0.135 \\
F1, validation (Sec.~\ref{sec:localisation}) & 0.185 & 0.260 \\
\midrule
F1, tampered area $<1$\,\%   & 0.040 & 0.123 \\
F1, tampered area 1--5\,\%   & 0.148 & 0.295 \\
F1, tampered area $>5$\,\%   & 0.333 & 0.413 \\
\midrule
F1, copy-move & 0.236 & 0.354 \\
F1, splicing  & 0.129 & 0.226 \\
\bottomrule
\end{tabular}
\end{table}

\subsection{Error Analysis}\label{sec:test-errors}
Under R, missed forgeries (126) were compared with detected ones (222), and false alarms (13) with correct
authentic verdicts (619); images in the uncertain band are excluded. The comparison uses automatically computed
image properties, with Fisher's exact test and a Holm correction~\cite{holm1979} over the 14 distinct tests
(Table~\ref{tab:test-fisher}; examples in Fig.~\ref{fig:failures}).

\begin{table}[t]
\centering
\caption{Missed vs detected forgeries (R): share of each property, ratio, Fisher $p$ and Holm-adjusted
$p$.}\label{tab:test-fisher}
\footnotesize
\setlength{\tabcolsep}{4pt}
\begin{tabular}{lccccc}
\toprule
Property & Missed & Det. & Ratio & $p$ & Holm $p$ \\
\midrule
Splicing                    & 60\,\% & 3\,\%  & \textbf{22$\times$}  & $<$0.001 & $<$0.001 \\
Keypoints, lowest tercile   & 49\,\% & 21\,\% & \textbf{2.4$\times$} & $<$0.001 & $<$0.001 \\
Tampered area $>5$\,\%      & 33\,\% & 57\,\% & 0.57$\times$ & $<$0.001 & $<$0.001 \\
JPEG source (vs TIFF)       & 44\,\% & 31\,\% & 1.4$\times$  & 0.015 & 0.16 \\
Low texture                 & 42\,\% & 34\,\% & 1.2$\times$  & 0.17  & n.s. \\
Flat, saturated, tiny area  & --     & --     & 1.1--1.4$\times$ & $\geq$0.28 & n.s. \\
\bottomrule
\end{tabular}
\end{table}

The causes of missed forgeries, in order of evidence, are as follows. (i)~\emph{Splicing}: the confident
``tampered'' verdicts are almost all copy-move images. Splicing leaves no internal duplicate, and after
resampling its compression and noise traces are weak (splicing vs authentic AUC 0.68 under R, 0.88 under B).
(ii)~\emph{Few keypoints}: the strongest module, copy-move keypoint matching, has little to work with in smooth,
low-detail images. (iii)~\emph{Area}: regions above 5\,\% are under-represented among misses (0.57$\times$), but
tiny regions ($<1$\,\%) are \emph{not} significantly over-represented. JPEG-sourced forgeries look
over-represented (1.4$\times$, $p=0.015$), but this does not survive the Holm correction; this is the residual
format effect first noted in Section~\ref{sec:classification}, and on test it even reverses under B, where
TIFF-sourced forgeries are detected less often (0.77 vs 0.89).

False alarms are only 13, and none of the image-property tags is significant after correction (low-keypoint
images: 0.33$\times$, $p=0.045$, Holm $p=0.45$). The copy-move features point clearly to the cause: \emph{all 13}
false alarms have copy-move keypoint matches that survive RANSAC~\cite{fischler1981ransac} (median 13 inliers), \emph{none} of the 619
correct authentic verdicts has any, and the false alarms also have many keypoints (median 1,500, the detector's
cap, against 732). Most of the eight false alarms in Fig.~\ref{fig:failures} are \emph{self-similar scenes}
(colonnades and pillars, repeated statues and ornaments, a lattice tower, a zebra-patterned sofa, rows of gaming
machines) in which genuinely repeated structure looks like a copy-move, a known weakness of copy-move
detection~\cite{christlein2012eval,wen2016coverage}. The scene descriptions are qualitative ($n=13$), but the inlier counts are
measured.

\begin{figure*}[t]
\centering
\includegraphics[width=0.9\textwidth]{fig_failures_R.png}
\caption{Test-set failures under protocol R. Row~1: most confident missed forgeries; row~2: random other missed
forgeries (red outline: true region); row~3: most confident false alarms; row~4: random other false alarms. Each
panel gives the forgery type and area class, the calibrated P(tampered) and a texture score. Most false alarms
are self-similar scenes (pillars, statues, a lattice tower, a zebra-patterned sofa, gaming machines).}
\label{fig:failures}
\end{figure*}
